\section{Related Work}
\label{sec:related}

\paragraph{Indirect prompt injection.}

Prompt injection inserts malicious text that causes an LLM to deviate from its
assigned objective \citep{perez2022ignore, willison2022pi}; in \emph{indirect} prompt
injection, the text is placed in third-party data that an application
retrieves \citep{greshake2023ipi}. The predominant
formulation of indirect prompt injection is \emph{instruction injection}: the
attacker embeds an instruction in external data, and the model follows it as
an instruction rather than treating it as data \citep{liu2024formalizing, yi2024bipia, debenedetti2024agentdojo}.

\paragraph{IPI benchmarks.}

IPI benchmarks measure whether adversarial content encountered during a
legitimate task can change an LLM application's output or an agent's behavior.
Benchmarks have moved from single-generation applications \citep{yi2024bipia}
and single tool calls \citep{zhan2024injecagent} to stateful, multi-step
environments with separate utility and attack checks
\citep{debenedetti2024agentdojo}, including dynamic tasks in which agents follow
benign third-party instructions \citep{li2026agentdyn}, and to web, computer-use,
and visual surfaces \citep{evtimov2025wasp, kuntz2025osharm, cao2026vpibench}.
Across these settings, evaluated payloads are mostly malicious instructions that
redirect the agent to attacker-chosen tasks; task information, if used, mainly
guides placement or coarse tailoring.

A more recent line couples injected content to the runtime decision it targets,
including attacks on tool selection, arguments, conditions, reasoning, and
delegation, as well as attacks propagated through memory, handoffs, skills,
and workflows \citep{wang2026agentpi, weng2026argus}. Once content is pertinent to
that decision, neither its relevance nor its wording separates an attack from
legitimate task evidence; what makes it an injection is that its source lacks
the authority, assigned by the system or user, to determine the decision. We
therefore make this authority explicit in every case: the user's request
designates which source determines the affected decision
(Section \ref{sec:threat}).

\paragraph{Defenses.}

Prompt-injection defenses differ in how they limit untrusted influence.
Prompting defenses such as Sandwich and Spotlighting
\citep{debenedetti2024agentdojo, hines2024spotlighting} and training methods such
as Instruction Hierarchy, StruQ, and SecAlign \citep{wallace2024instruction, chen2024struq, chen2025secalign} shape how the model treats trusted
instructions versus external data. These methods rely on the model to apply
the intended boundary and do not independently constrain the calls it emits.

Detection-based defenses instead classify or remove suspected injections
before they influence the agent \citep{li2025piguard, shi2025promptarmor}.
Their coverage therefore depends on the injection being recognizable from
the content that carries it or the wider observed trajectory.

Plan- and policy-based defenses constrain proposed actions against a
task-derived action space. AgentDojo's Tool Filter \citep{debenedetti2024agentdojo},
Progent \citep{shi2025progent}, and DRIFT \citep{li2025drift} check calls against a tool
set, policy, or plan, respectively; DRIFT also uses an LLM-based isolator to
mask tool-output instructions that conflict with the user request. These
systems enforce conformity to the represented plan or policy, whose coverage
depends on how completely that policy captures the task. The policy or plan generally does
not constrain the content of free-text arguments.

Architectural isolation and information-flow control constrain how content can
reach agent components and actions. IsolateGPT and ACE
\citep{wu2025isolategpt, li2025ace} isolate applications or execution steps, while
CaMeL and FIDES \citep{debenedetti2025defeating, costa2025securing} track
provenance labels on values and enforce policies at tool calls. These systems
enforce component, source, and flow policies; they do not by themselves
establish the truth or topic-specific authority of content admitted by those
policies.

Evidence- and authorization-based defenses check whether benign runtime
evidence entails an action, as in ARGUS \citep{weng2026argus}, or whether its
parameters come from sources authorized by the request, as in AuthGraph
\citep{authgraph2026}; ARGUS still relies on judging which content is benign, and
AuthGraph on the authorization graph it derives from the request.
